\emph{Tree stabiliser on \(K\) with a tree cubic term.} Read the cubic term of Equation~\eqref{eq:Smatrix} on the flux directions of \(K\), specialised there to \(C_{ijk}=\mathrm i c_3\epsilon_{ijk}/g^2\) on the flux directions with constant \(c_3\) (a coupling label local to these two paragraphs, not the quartic-potential \(c_3\) of the \(U(4)\) discussion), and write \(V_{\mathrm t}=\bar a(\chi^4-\tfrac23c_3\chi^3)\), \(r_3=c_3/2\chi\), \(\kappa_3=2-c_3/\chi=2(1-r_3)\). The tree part is stationary at \(\chi=c_3/2\) and convex there, \(\partial^2_{\ln\chi}(\rho_MV_{\mathrm t}/g^2)=4\tilde U(4-3r_3)>0\). At \(X^i=\chi J^i\) on \(S^2_N\), with \(J^i\) the spin-\((N-1)/2\) generators, the difference of the quadratic forms of the quartic and cubic terms at \(c_3\) and at \(c_3=0\) is exactly \(\tfrac12(c_3/\chi)\,\chi^2\) times the real symmetric operator \((\mathsf{SO}\,A)^i=-\epsilon_{kij}\mathrm i[J^k,A^j]\); this is an identity between two unfixed quadratic forms, checked at \(N\le5\) with residual below \(1.5\times10^{-13}\). In the same reduction of the fluctuation operator that gives the vector spectrum of the preceding paragraph, the cubic term therefore changes only the spin-orbit coefficient, and the vector spectrum becomes \(\chi^2(l(l+1)-\kappa_3l)\), \(\chi^2(l(l+1)+\kappa_3)\), \(\chi^2(l(l+1)+\kappa_3(l+1))\) with multiplicities \(2l+3\), \(2l+1\), \(2l-1\) over \(l=1,\dots,N-1\), of total \(3(N^2-1)\); at \(\kappa_3=2\) it is the vector spectrum of the preceding paragraph, and the spinor and ghost parts being unchanged, there is no negative mode exactly for \(-1\le\kappa_3\le2\) (at \(\kappa_3=-2\) the mode \(-2\chi^2\) appears). The supplied example of the preceding paragraph lies on this branch with \(c_3/\chi\simeq2.34\) and \(\kappa_3\simeq-0.336\), inside the window.

\emph{The tree curve of the cubic branch.} Write \(U_{\chi\chi}=4\tilde U(4-3r_3+\delta)\) at a stationary point of the branch, which defines \(\delta\). With \(\partial_{\ln\chi}(\rho_MV_{\mathrm t}/g^2)=4\tilde U(1-r_3)\), the value \(\beta_K=-2\partial_{\ln\chi}(\rho_MV_{\mathrm t}/g^2)/U_{\chi\chi}\) recorded above reads \(\beta_K=2(r_3-1)/(4-3r_3+\delta)\) on the branch; \(\beta_K=1\) is equivalent to \(r_3=(6+\delta)/5\), and at \(\kappa_3=-2\) one has \(\beta_K=2/(\delta-2)\). In the units of the example of the preceding paragraph its supplied values give \(r_3\simeq1.168\), \(\kappa_3\simeq-0.336\), \(\delta\simeq-0.159\) and \(\beta_K\simeq2\cdot5.34/10.68\simeq1.000\), a restatement of the value quoted there. If in addition the large-\(N\) weights on this branch give \(\hat V/\hat T\to\kappa_3^2/4\) (supplied, not recomputed here), then at a stationary point \(\hat V/\hat T=(1-r_3)^2\), so \(\hat V=\hat T\) requires \(r_3=0\) or \(r_3=2\): the first is \(c_3=0\), the second leaves the window of no negative modes. On that input no stationary point of the cubic branch has \(\hat V=\hat T\) with no negative mode, and the example has \(\hat V/\hat T\simeq0.028\). The one-loop curvature \(\delta\), the weight ratio, and the response of the one-loop density beyond \(\sqrt{|g|}\) are not computed here.
